\documentclass{stud-conf-pmk}

\usepackage{multirow}
\usepackage{graphicx}

\begin{document}
\year{2026}

\begin{paper}
\author[Антонов М.\,Е.]{Антонов Михаил Евгеньевич}
\title{Мультизадачный анализ музыкальных композиций с~помощью предобученного аудио-трансформера}
\speciality{Прикладная математика и информатика}
\student{Студент магистратуры}
\presented{Представил доцент кафедры информационных технологий М.\,Ю.\,Кудряшов}
\UDC{004.934, 004.85}
\abstract{Реализована мультизадачная модель на~основе предобученного Audio Spectrogram Transformer с~общим стволом и~четырьмя классификационными головами: сегментация структуры, классификация возбуждения и~валентности, определение жанра. Обучение проведено на~датасетах DEAM, Harmonix~Set и~GTZAN методом Round-Robin.}
\keywords{глубокое обучение, мультизадачное обучение, Audio Spectrogram Transformer, музыкальный информационный поиск, классификация аудио}
\maketitle

\introduction

Музыкальный информационный поиск (MIR) объединяет ряд задач: определение жанра трека, оценку его эмоционального содержания, выделение структурных секций (вступление, куплет, припев и~т.\,д.). На~практике каждая из~них решается отдельной моделью, что~требует независимого хранения весов и~игнорирует общую акустическую природу входных данных. Парадигма мультизадачного обучения~\cite{bib_ruder_mtl} позволяет обойти это ограничение за~счёт совместного использования общего ствола сети.

В~области аудиоанализа стандартом сегодня являются модели на~основе механизма самовнимания~\cite{bib_vaswani_transformer}; одной из~наиболее успешных среди них является Audio Spectrogram Transformer (AST)~\cite{bib_gong_ast}~--- адаптация Vision Transformer для~спектрограммного представления звука, предобученная на~датасете AudioSet (около 2~млн размеченных аудиоклипов). В~настоящей работе AST используется как~общий ствол мультизадачной модели, поверх которого размещены четыре классификационные головы.


\section{Архитектура модели}

Входной аудиосигнал длительностью 10\,с с~частотой дискретизации $22{,}05$\,кГц преобразуется в~лог-мел-спектрограмму размера $128 \times T$, где 128~--- количество мел-фильтров, $T$~--- число временных кадров. Спектрограмма разбивается на~патчи $16 \times 16$, проецируемые линейным слоем в~768-мерное пространство; к~полученной последовательности добавляются обучаемый служебный токен \textsc{cls} и~позиционные эмбеддинги.

Последовательность обрабатывается стеком из~12~блоков трансформера; каждый блок содержит многоголовое самовнимание (12~голов) и~полносвязную сеть с~расширением до~3072 нейронов. Веса первых 8~блоков инициализируются из~AST, предобученного на~AudioSet, и~остаются замороженными. Последние 4~блока дообучаются на~целевых данных, что~помогает избежать катастрофического забывания признаков, выученных на~AudioSet. Применяется дифференциальная скорость обучения: $5 \cdot 10^{-5}$ для~бэкбона и~$10^{-3}$ для~голов.

С~выхода последнего блока берётся эмбеддинг \textsc{cls}-токена ($\mathbb{R}^{768}$), агрегирующий информацию о~всём фрагменте, и~подаётся в~общий полносвязный слой $768 \to 256$; затем результат поступает на~четыре независимые головы. Общее число параметров~--- 87\,млн, из~них 14\,млн обучаемые (4~блока трансформера~+ общий слой~+ головы). Схема приведена на~рис.~\ref{fig:ast}.

\begin{figure}
\centering
\includegraphics[height=0.55\textheight]{figures/ast_architecture.pdf}
\caption{Архитектура мультизадачной модели на~основе AST.}
\label{fig:ast}
\end{figure}


\section{Мультизадачное обучение}

Каждая из~четырёх голов~--- линейный слой $256 \to K_i$ с~последующей функцией \textsc{softmax}; число классов $K_i$ для~каждой задачи приведено в~табл.~\ref{tab:heads}.

\begin{table}
\centering
\caption{Классификационные головы модели.}
\begin{tabular}{|l|c|l|}
\hline
Задача & $K$ & Классы \\
\hline
Сегментация структуры & 6 & Intro, Verse, Bridge, Chorus, Instrumental, Outro \\
\hline
Возбуждение (arousal) & 3 & low, mid, high \\
\hline
Валентность (valence) & 3 & negative, neutral, positive \\
\hline
Жанр & 10 & blues, classical, country, disco, hiphop, \\
& & jazz, metal, pop, reggae, rock \\
\hline
\end{tabular}
\label{tab:heads}
\end{table}

Частная функция потерь $\mathcal{L}_i$ для~каждой задачи~--- взвешенная кросс-энтропия со~сглаживанием меток ($\varepsilon = 0{,}1$); веса классов обратно пропорциональны их~частотам в~обучающей выборке, что~компенсирует дисбаланс (особенно выраженный для~сегментации, где Verse и~Chorus встречаются в~3--5~раз чаще Bridge и~Outro). Совокупная функция потерь:
$$\mathcal{L} = \sum_{i=1}^{4} \lambda_i \mathcal{L}_i,$$
где $\lambda_{\text{seg}} = 1{,}5$, $\lambda_i = 1{,}0$ для~остальных задач (значения подобраны по~валидационной выборке).

Используемые датасеты содержат метки только для~одной из~задач: DEAM~--- эмоции, Harmonix~Set~--- структура, GTZAN~--- жанр. Поэтому применяется стратегия Round-Robin: батч на~каждом шаге обучения формируется только из~одного датасета, потеря вычисляется только по~активным головам. Общий ствол модели при~этом поочерёдно получает градиенты от~всех задач, реализуя жёсткое разделение параметров (\textit{hard parameter sharing}).


\section{Эксперименты и~результаты}

Обучающая выборка составлена из~трёх открытых датасетов: DEAM (1802~трека с~непрерывными аннотациями возбуждения и~валентности), Harmonix~Set (743~трека с~разметкой структуры) и~GTZAN (999~треков с~разметкой по~10~жанрам). Совокупный объём данных~--- 3544~трека. Окна длительностью 10~секунд извлекаются с~шагом 1~секунда, что~даёт около 176\,000 обучающих примеров.

Обучение проводилось в~течение 8~эпох с~использованием оптимизатора AdamW и~планировщика \textsc{CosineAnnealingWarmRestarts}. Размер эффективного батча~--- 64 (16~примеров с~накоплением градиента за~4~шага). Применялась аугментация SpecAugment~--- маскирование случайных частотных и~временных полос лог-мел-спектрограммы.

Результаты на~тестовой выборке приведены в~табл.~\ref{tab:results}.

\begin{table}
\centering
\caption{Метрики качества предложенной модели.}
\begin{tabular}{|l|c|c|c|}
\hline
Задача & Accuracy & Macro-F1 & Случайная базовая \\
\hline
Сегментация (6 классов) & 31{,}4\,\% & 31{,}5\,\% & 16{,}7\,\% \\
\hline
Возбуждение (3 класса) & 65{,}9\,\% & 65{,}3\,\% & 33{,}3\,\% \\
\hline
Валентность (3 класса) & 55{,}9\,\% & 56{,}0\,\% & 33{,}3\,\% \\
\hline
Жанр (10 классов) & 82{,}7\,\% & 81{,}9\,\% & 10{,}0\,\% \\
\hline
\end{tabular}
\label{tab:results}
\end{table}

На~задаче определения жанра модель выходит на~уровень, сопоставимый с~однозадачными моделями для~GTZAN (82{,}7\,\% против~85--90\,\% у~SOTA). Возбуждение и~валентность~--- 65{,}9\,\% и~55{,}9\,\% соответственно при~базовой линии 33{,}3\,\%; разница между~ними отражает известную в~литературе большую субъективность разметки валентности. Сегментация (31{,}4\,\% при~базовой 16{,}7\,\%) остаётся наиболее сложной задачей: помимо дисбаланса классов, на~качество влияет неоднозначность ручной разметки границ секций (расхождения между~аннотаторами в~Harmonix~Set достигают $\pm 3$\,с).


\conclusion

Реализована мультизадачная модель на~основе предобученного AST с~общим стволом и~четырьмя классификационными головами. Round-Robin-обучение на~трёх датасетах с~непересекающимися наборами меток позволило получить рабочее качество по~всем четырём задачам в~рамках одной модели в~87\,млн параметров: 82{,}7\,\% для~жанра, 65{,}9\,\% и~55{,}9\,\% для~возбуждения и~валентности, 31{,}4\,\% F1 для~сегментации.

Дальнейшие шаги: добавление регрессионных голов для~непрерывных шкал возбуждения и~валентности; постобработка предсказаний сегмента (декодирование Витерби с~позиционными априорными вероятностями); расширение набора задач~--- определение темпа и~тональности.


\begin{thebibliography}{9}

\bibitem{bib_gong_ast}
Gong, Y. AST: Audio Spectrogram Transformer~/ Y.\,Gong, Y.-A.\,Chung, J.\,Glass~// Proceedings of Interspeech.~--- 2021.~--- P.\,571--575.~--- DOI: \url{https://doi.org/10.21437/Interspeech.2021-698}

\bibitem{bib_ruder_mtl}
Ruder, S. An Overview of Multi-Task Learning in Deep Neural Networks~// arXiv preprint.~--- 2017.~--- arXiv:1706.05098.~--- URL: \url{https://arxiv.org/abs/1706.05098}

\bibitem{bib_vaswani_transformer}
Vaswani, A. Attention Is All You Need~/ A.\,Vaswani, N.\,Shazeer, N.\,Parmar [et~al.]~// Advances in Neural Information Processing Systems (NeurIPS).~--- 2017.~--- Vol.\,30.~--- P.\,5998--6008.

\end{thebibliography}

\end{paper}
\end{document}
